\documentclass{article}

\usepackage{neurips_2026}

\usepackage[utf8]{inputenc}
\usepackage[T1]{fontenc}
\usepackage{hyperref}
\usepackage{url}
\usepackage{booktabs}
\usepackage{amsfonts}
\usepackage{amsmath}
\usepackage{nicefrac}
\usepackage{microtype}
\usepackage{xcolor}
\usepackage{algorithm}
\usepackage{algpseudocode}
\usepackage{multirow}
\usepackage{enumitem}

\title{DPGrammar: Joint Viterbi Repair for Grammar-Constrained\\Diffusion Language Models}

\author{%
  Anonymous Author(s)\\
  Affiliation\\
  Address\\
  \texttt{email}
}

\begin{document}

\maketitle

\begin{abstract}
Incremental parsers are left-to-right; masked diffusion language models are not.
A diffusion decoder commits tokens in confidence order, so a grammar can only be
enforced at the parse frontier, and every token placed beyond it is an
unvalidated fact that the parser may later reject. We present \textbf{DPGrammar},
a repair layer that treats such a rejection as a joint decision over the whole
violated span rather than a single-token fix. The repair is a Viterbi search over
(parser state $\times$ position) whose candidate set is drawn \emph{from the
automaton}---the tokens the parser can actually consume---rather than from the
model's top-$k$ over the vocabulary. This distinction is not cosmetic: on
JSONSchemaBench \texttt{medium}, $46\%$ of live parser states admit fewer than
$100$ of $126{,}349$ tokens, so ranking the vocabulary and filtering afterwards
is rejection sampling against a small target and fails outright on $49.3\%$ of
repairs. Drawing from the automaton makes failure structurally impossible and
takes the dead-end rate to $0\%$.
On $n{=}511$ instances, DPGrammar reaches $96.7\%$ schema validity against
$91.8\%$ for the same system without the DP layer and $78.3\%$ for LAVE, the
strongest published diffusion baseline, while collapsing mean repair events from
$33.2$ to $1.2$.
We further show that schema validity alone is a misleading objective: it scores
\texttt{\{"objects": []\}} as a success and an $851$-character document that does
not quite close as a failure. Pairing validity with a content floor reverses the
ranking between the two baselines and leaves DPGrammar first at every threshold.
\end{abstract}

\section{Introduction}

Masked diffusion language models~\citep{llada,dream} generate by iteratively
unmasking positions in confidence order rather than left to right. This is what
makes them fast, and it is also what makes grammar-constrained decoding hard.
An incremental parser---the standard machinery for constrained
decoding~\citep{outlines,dong2024xgrammar,microsoft2025llguidance}---consumes
one token at a time from the left. At any point during diffusion the sequence
looks like
\[
\underbrace{\texttt{[consumed prefix]}}_{\text{parser has validated}}
\;\underbrace{\texttt{[frontier]}}_{\text{constrainable}}
\;\texttt{[MASK]}\;\underbrace{\texttt{[token]}}_{\text{unvalidated}}\;\texttt{[MASK]}\;\cdots
\]
The grammar can be enforced only at the frontier. Everything the model has
already placed beyond it was sampled without constraint, and---because the
sampler is monotone, $x_0 \leftarrow \mathbf{1}[\text{mask}]\odot x_0 + (1-\mathbf{1}[\text{mask}])\odot x$---it
is already final. When the frontier advances into those positions, the parser may
reject, and the decoder must repair a commitment it cannot take back.

Existing diffusion decoders handle this in one of three ways. They force
near-autoregressive ordering, giving up the parallelism that motivates diffusion;
they verify each frontier token by sampling $K$ random completions and abort a
whole block on failure~\citep{zhang2026lookaheadthenverifyreliableconstraineddecoding};
or they remask the offending position and let the sampler try again. The third is
cheap but myopic: it repairs one token at a time and cannot see that the
violation is joint across several positions.

\paragraph{Contributions.}
\begin{enumerate}[leftmargin=1.4em,itemsep=1pt,topsep=2pt]
\item \textbf{DPGrammar}, a joint repair layer that solves the violated span with
      a Viterbi search over parser states (\S\ref{sec:method}). On $n{=}511$
      JSONSchemaBench \texttt{medium} instances it reaches $96.7\%$ schema
      validity versus $91.8\%$ without the layer, and reduces mean repair events
      from $33.2$ to $1.2$.
\item \textbf{Candidates from the automaton, not from the vocabulary}
      (\S\ref{sec:cand}). We show that proposing the model's top-$k$ and
      filtering by the grammar is rejection sampling against a target set that is
      usually tiny---$46\%$ of live states admit fewer than $100$ of $126{,}349$
      tokens---and that drawing from the parser's outgoing edge set instead makes
      repair failure structurally impossible ($49.3\%\!\to\!0\%$ dead ends).
\item \textbf{A grammar-derived termination rule} (\S\ref{sec:stop}) that
      replaces a step-count threshold with the question ``can the parser do
      anything other than stop?'', unifying JSON and SMILES under one rule with
      no per-dataset constant.
\item \textbf{A measurement correction} (\S\ref{sec:metric}). The completion flag
      commonly logged as \texttt{valid} is not schema validity; it overstates
      grammar-constrained methods by $4.6$--$6.7$\,pp and unconstrained ones by
      $0$\,pp. We also show that validity alone rewards degenerate output, and
      report it against a content floor.
\end{enumerate}

\section{Related Work}

\paragraph{Constrained decoding for diffusion.}
\citet{cd4dllm} propose the first CFG-constrained dLLM decoder. DINGO
\citep{suresh2025dingoconstrainedinferencediffusion} applies dynamic programming
for syntactic correctness at a higher per-step cost, over the full sequence
rather than over a violated span. LAVE
\citep{zhang2026lookaheadthenverifyreliableconstraineddecoding} verifies each
frontier token by sampling $K$ random completions, accepting only if one is
grammar-valid; it is the dominant published baseline. Plan-style decoders
\citep{li2026planverifyfillstructured} emphasise structure-aware planning. We
stay at the token level with an explicit incremental matcher, and add joint
search only where the parse actually conflicts.

\paragraph{Constrained decoding for autoregressive models.}
Outlines~\citep{outlines}, XGrammar~\citep{dong2024xgrammar} and
llguidance~\citep{microsoft2025llguidance} compile grammars into incremental
automata and mask logits per step. All assume left-to-right commitment; the
question of what to do when a token is committed \emph{out of order} does not
arise, and it is precisely the question diffusion raises.

\paragraph{Speed.}
Speculative decoding~\citep{leviathan2023fastinferencetransformersspeculative}
and sparse automaton acceleration~\citep{su2026vectorizingtrieefficientconstrained}
reduce the cost of verification. Our repair layer is complementary: it fires on
roughly one position per instance and is not on the critical path of a clean
step.

\section{Problem Setting}
\label{sec:problem}

Let $G$ be a grammar with an incremental matcher exposing four operations:
\texttt{try\_consume}, \texttt{rollback}, \texttt{clone}, and a mask
$\mathcal{L}(q)\subseteq V$ of tokens consumable from state $q$. Let the decoder
hold a sequence $x$ over $V\cup\{\textsc{mask}\}$ and a frontier index $c$ such
that $x_{<c}$ has been consumed, leaving the matcher in state $q_c$.

At each denoising step the model proposes tokens at masked positions. Frontier
masking guarantees $x_c\in\mathcal{L}(q_c)$ when $x_c$ is chosen this step, but
positions $>c$ are unconstrained. After placement the decoder advances the parser
over the newly contiguous region; if it stops at position $v$ with
$x_v\notin\mathcal{L}(q_v)$, we call $v$ a \emph{violator}.

\paragraph{Repair as search.}
The repair problem is: find $y$ over a span $[v,e)$ maximising
$\sum_{p} \log p_\theta(y_p\mid x)$ subject to the matcher accepting
$y_v\cdots y_{e-1}$ from $q_v$. This is a shortest-path problem on a layered
graph whose nodes are (position, parser state) and whose edges out of a state are
exactly $\mathcal{L}(q)$. Two design choices determine whether the search is
sound, and both were wrong in the natural implementation.

\section{Method}
\label{sec:method}

\subsection{The repair cascade}

When the parser rejects at $v$, the decoder tries, in order:

\begin{enumerate}[leftmargin=1.4em,itemsep=1pt,topsep=2pt]
\item \textbf{Termination check.} If the parser cannot consume anything, or the
      model has placed a stop token at the frontier and the parser permits
      stopping, the document ends (\S\ref{sec:stop}).
\item \textbf{Single-token retry.} Walk the ranked logits at $v$ for up to $10$
      candidates and take the first the parser accepts. This resolves $62\%$ of
      violations at negligible cost.
\item \textbf{Joint DP repair.} Otherwise solve the span $[v,e)$ jointly
      (\S\ref{sec:dp}).
\item \textbf{Hand back.} If no assignment exists, remask $v$ and let the sampler
      re-derive it.
\end{enumerate}

The span end $e$ is $\min(v+L,\,\text{first \textsc{mask}})$ with $L$ a compute
budget; the DP does not fill holes, only rewrites committed positions.

\subsection{Joint repair by Viterbi over parser states}
\label{sec:dp}

We maintain a map from parser state to the best path reaching it:
\[
S_p:\; q \;\mapsto\; (\text{matcher clone},\; \text{edits},\; \textstyle\sum\log p_\theta).
\]
Two paths that reach the same state are merged, keeping the higher-scoring one;
this is what makes the search a dynamic program rather than an enumeration, and
it is sound because the state determines the future. At the final layer we take
the highest-scoring state and backtrack.

The state key is the matcher's own mask, $\texttt{bytes}(\mathcal{L}(q))$, which
the implementation already computes; distinct keys are distinct automaton nodes.
In practice $|S_p|$ stays in the single digits, so the cost is
$O(\text{span}\times|S|\times k)$ matcher operations with $|S|$ bounded by the
reachable automaton rather than by $k^{\text{span}}$.

\subsection{Candidates from the automaton}
\label{sec:cand}

The natural implementation asks the model for its top-$k$ tokens and then asks
the parser which of them it accepts. This is rejection sampling against
$\mathcal{L}(q)$, and it fails whenever $\mathcal{L}(q)$ is small relative to
$|V|$. Measuring $1{,}687$ live states on JSONSchemaBench \texttt{medium} with
$|V|=126{,}349$, the distribution of $|\mathcal{L}(q)|$ is sharply bimodal:
\[
p_{25}=94,\quad p_{50}=101,\quad p_{75}=125{,}311,
\]
with $46\%$ of states admitting fewer than $100$ tokens. At those states---key
boundaries, enumerations, closings---a top-$100$ draw from a $126$k vocabulary
frequently misses $\mathcal{L}(q)$ entirely, and the DP reports ``no path'' where
a path exists. At a freshly opened JSON object we measured $|\mathcal{L}(q)|=2$.

We therefore invert the order. For each live state we read $\mathcal{L}(q)$ and
take the $k$ most probable tokens \emph{within it}:
\begin{equation}
\mathcal{C}(q,p) \;=\; \operatorname*{top-}k_{\,t\,\in\,\mathcal{L}(q)\setminus\{\textsc{eos}\}}\; \log p_\theta(t\mid x,p).
\end{equation}
Since a live state always admits at least one token---$0$ of $2{,}028$ sampled
states had an empty mask---the candidate set is never empty and the DP can always
complete the span. Where $|\mathcal{L}(q)|\le k$ the edge set is enumerated
exactly rather than sampled; this holds at $67\%$ of state-positions. Empirically
the dead-end rate falls from $49.3\%$ to $0\%$, repair events fall by $46\%$, and
wall time falls slightly. The cost is one mask read per (state, position), which
the implementation already paid to compute state keys.

\paragraph{Why $\textsc{eos}$ is excluded.}
The stop token is marked legal by the mask at accepting states but is not a
grammar symbol: \texttt{try\_consume} refuses it, because there is nothing to
advance over. Offering it wastes a candidate slot and, where it is the only
``legal'' token, empties the set and imitates a dead end.

\subsection{Termination as a grammar question}
\label{sec:stop}

Because \textsc{eos} is not consumable, a finished document arrives at the
frontier looking like a violation, and something must recognise it. The natural
test---``is the parser in an accepting state?''---is wrong, and the error is
invisible on JSON. It asks whether the parser \emph{may} stop, not whether it
\emph{must}. On $53$ of $53$ completed JSON documents the accepting state admitted
exactly one token, \textsc{eos}, so the two questions coincide. On SMILES a
complete molecule such as \texttt{CCO} is accepting with $1{,}388$ tokens still
legal, and stopping there takes the length decision away from the model: every
molecule came out as \texttt{C} or \texttt{CC}.

We use two rules, for two different decisions:
\begin{itemize}[leftmargin=1.4em,itemsep=1pt,topsep=2pt]
\item \emph{The harness ends the document} only when
      $\mathcal{L}(q)\setminus\{\textsc{eos},\textsc{eot}\}=\emptyset$ and $q$ is
      accepting---the parser has no alternative, so ending is not a choice
      anyone is making.
\item \emph{The model ends the document} when it places a stop token at the
      frontier and $q$ is accepting---the model has asked, and the grammar need
      only permit it.
\end{itemize}
This replaces a step-count threshold with a property of the grammar. On JSON it
reproduces the previous behaviour exactly; on SMILES it recovers median output
length $3\!\to\!8$ characters and RDKit validity $62.5\%\!\to\!81.2\%$ while
eliminating all repair activity ($16$ DP calls and $16$ remasks $\to$ $0$).

\subsection{System optimisations}

The layer sits on a frontier-constrained decoder whose engineering we treat as
infrastructure rather than contribution: an incremental matcher carried across
denoising steps instead of re-parsing the prefix, additive-increase batching of
committed tokens, in-place resampling that reuses the current logits, and
overlapping mask construction with the next forward pass. Throughout the paper,
``\textbf{no-DP baseline}'' denotes exactly this system with the repair layer
disabled, so that every comparison isolates the DP layer.

\section{Evaluation Methodology}
\label{sec:metric}

\subsection{The completion flag is not validity}

Drivers in this setting commonly log a boolean \texttt{valid} taken from the
generator, defined as ``a stop token is present and no position before it is
masked''. That is \emph{completion}, not validity: an unconstrained decoder can
finish cleanly on malformed JSON and score a success. Recomputing schema validity
with \texttt{jsonschema} over the stored schemas gives

\begin{center}
\begin{tabular}{lcc}
\toprule
 & completion flag & schema validity \\
\midrule
LAVE                 & $78.3\%$ & $78.3\%$ \\
No-DP baseline       & $98.4\%$ & $91.8\%$ \\
DPGrammar            & $96.7\%$ & $96.7\%$ \\
\bottomrule
\end{tabular}
\end{center}

The bias is not uniform. Grammar-constrained methods emit documents that
terminate cleanly but are structurally broken---one closing brace too many, a key
name written over a value---which count as complete but not valid. LAVE falls
back to autoregressive decoding on failure and produces well-formed JSON, so for
it the two notions coincide. Any comparison that mixes the two measures the
fallback strategy rather than the constraint algorithm. All numbers we report are
recomputed schema validity.

\subsection{Validity alone rewards giving up}

Schema validity is monotone in nothing that a reader cares about. On instance
\texttt{o17539} the no-DP baseline emits \texttt{\{"objects": []\}}---zero
leaves, schema-valid, scored a win---while DPGrammar emits $851$ characters of
populated objects that do not quite close, scored a loss. On \texttt{o11975} the
baseline's ``success'' is
\texttt{\{"faction":"rebel","pilots":[\{"name":"0","ship":"0"\}],"version":"0.0.0"\}}.
Across the benchmark, $22\%$ of the baseline's valid outputs carry fewer than
five leaves.

We therefore report validity against a content floor: an output counts only if it
validates \emph{and} carries at least $\kappa$ scalar leaves. We report
$\kappa\in\{1,3,5,10,15\}$ so that the threshold is visibly not doing the work.

\subsection{Setup}

LLaDA-8B-Instruct~\citep{llada}, $T{=}128$ denoising steps, generation length
$256$, block length $32$, temperature $0.2$, on a single A100 per shard.
Grammars are compiled by llguidance~\citep{microsoft2025llguidance}.
JSONSchemaBench~\citep{geng2025jsonschemabenchrigorousbenchmarkstructured}
\texttt{medium} provides $586$ instances; $75$ have schemas llguidance cannot
compile, so no method runs on them and they are excluded from every arm, leaving
$n{=}511$. LAVE is reimplemented on the same backbone and step budget with a
$120$\,s per-instance timeout.

\section{Results}

\begin{table}[t]
\centering
\small
\caption{JSONSchemaBench \texttt{medium}, $n{=}511$. Validity is recomputed with
\texttt{jsonschema}; ``$\kappa$'' columns additionally require $\kappa$ scalar
leaves. Repair events are remasks for the grammar-based methods and lookahead
draws for LAVE.}
\label{tab:main}
\begin{tabular}{lccc}
\toprule
 & LAVE & No-DP baseline & \textbf{DPGrammar} \\
\midrule
Schema validity              & $78.3\%$ & $91.8\%$ & $\mathbf{96.7\%}$ \\
\;\;valid $\wedge\ \kappa{=}1$   & $78.3\%$ & $84.9\%$ & $\mathbf{93.7\%}$ \\
\;\;valid $\wedge\ \kappa{=}3$   & $76.7\%$ & $79.1\%$ & $\mathbf{90.6\%}$ \\
\;\;valid $\wedge\ \kappa{=}5$   & $74.0\%$ & $74.6\%$ & $\mathbf{84.9\%}$ \\
\;\;valid $\wedge\ \kappa{=}10$  & $56.0\%$ & $50.7\%$ & $\mathbf{60.5\%}$ \\
\;\;valid $\wedge\ \kappa{=}15$  & $33.7\%$ & $29.5\%$ & $\mathbf{35.0\%}$ \\
\midrule
Median leaves (valid only)   & $13$ & $11$ & $12$ \\
Median chars (valid only)    & $530$ & $385$ & $466$ \\
\midrule
Mean repair events           & $253.46$ & $33.23$ & $\mathbf{1.19}$ \\
Timeouts ($>$120\,s)         & $60$ & $0$ & $0$ \\
\midrule
Mean time (s)                & $38.50$ & $\mathbf{13.45}$ & $16.75$ \\
Median time (s)              & $25.70$ & $\mathbf{13.40}$ & $15.05$ \\
P95 time (s)                 & $120.00$ & $\mathbf{21.43}$ & $36.21$ \\
\bottomrule
\end{tabular}
\end{table}

\paragraph{The DP layer is what buys content, not only validity.}
Table~\ref{tab:main} shows the no-DP baseline ahead of LAVE on raw validity by
$13.5$\,pp. The lead shrinks monotonically as the content floor rises---to
$6.6$\,pp at $\kappa{=}1$, $2.4$\,pp at $\kappa{=}3$, and $0.6$\,pp at
$\kappa{=}5$---and reverses beyond it, to $-5.3$\,pp at $\kappa{=}10$ and
$-4.2$\,pp at $\kappa{=}15$. Median length of its valid outputs is $385$
characters against LAVE's $530$. A substantial part of the baseline's apparent
advantage is bought by emitting the smallest document the schema permits, and
the advantage disappears exactly when the metric stops accepting such
documents.

DPGrammar does not make that trade. It is first at every threshold, and its valid
outputs are longer and richer than the baseline's ($466$ characters, $12$ leaves
versus $385$ and $11$). This is the result the raw-validity table conceals: the
DP layer's contribution is not a few more valid documents, it is valid documents
that still say something.

\paragraph{Cost.}
Mean repair events fall from $33.2$ to $1.19$, a $28\times$ reduction, for a
$3.3$\,s increase in mean wall time. LAVE's stochastic verification costs
$253.5$ draws per instance and times out on $60$ of $511$; both grammar-based
methods time out on none.

\section{Analysis}

\subsection{What the repair layer actually changes}

We instrumented every repair site. Of $118$ violations on a $64$-instance
subsample, $62\%$ are resolved by the single-token retry and never reach the DP.
Of those that do, the automaton candidate set eliminates dead ends entirely, and
the DP writes a mean of $5.95$ tokens over a mean span of $14.0$.

\subsection{Ablations, including four that changed nothing}
\label{sec:ablation}

We report these because they delimit where the gain does and does not come from.

\begin{center}
\begin{tabular}{llc}
\toprule
Component & Varied & Effect on validity \\
\midrule
Candidate source     & vocabulary top-$k$ $\to$ automaton & dead ends $49.3\%\to0\%$ \\
Repair window        & single token $\to$ full span       & 1-char leaves $22.3\%\to7.3\%$ \\
Termination rule     & accepting $\to$ nothing-consumable & SMILES $62.5\%\to81.2\%$ \\
\midrule
Objective            & $\log p$ $\to$ min-edit            & none ($-0.8$ edits) \\
Candidate admission  & always admit committed token       & none (identical output) \\
Post-hoc enrichment  & on $\to$ off                       & none (identical output) \\
Edit penalty         & $\lambda\in\{0,2\}$                & none (never non-zero) \\
\bottomrule
\end{tabular}
\end{center}

\paragraph{The objective does not matter; the search space does.}
Replacing $\log p$ with a lexicographic minimum-edit objective reduces mean edits
per repair only from $5.95$ to $5.15$. Since minimum-edit attains the minimum by
construction, $87\%$ of the rewriting the DP performs is \emph{forced} by the
grammar given the violation, not chosen by the objective. Changing what the
search optimises therefore cannot help; changing what it searches over---the
candidate set---moves the dead-end rate from $49.3\%$ to $0$.

\subsection{What repair cannot fix}

llguidance compiles a JSON schema into a grammar that fixes the order of an
object's properties to their declaration order: \texttt{\{"a":1,"b":2\}} is
rejected when the schema declares \texttt{b} first, although both orders satisfy
the schema and \texttt{jsonschema} accepts either. A model that writes the
properties in a different order therefore produces a violation that is not an
error, and the correct repair---moving a block---is not expressible by a
positional DP, which can only overwrite in place. The repair then writes the
demanded key name over a value intended for another field. We measured this class
at $8\%$ of violations; at states where the grammar admits exactly one token,
$56\%$ of violations are of this kind.

\section{Limitations}

\paragraph{Metric threshold.}
The content floor $\kappa$ is a reporting choice, not a learned quantity. We
report five values and the ranking of DPGrammar is stable across all of them, but
the two baselines exchange places between $\kappa{=}3$ and $\kappa{=}5$, so any
single-threshold claim about them would be fragile.

\paragraph{Scope.}
Results are one backbone (LLaDA-8B-Instruct), one benchmark split, and one seed.
The SMILES evidence in \S\ref{sec:stop} is a $16$-instance pilot and establishes
only that the termination rule generalises across grammar formalisms, not that
the validity gain does.

\paragraph{Remaining failures.}
Of the residual failures we inspected, all were grammar-valid prefixes: the
documents were incomplete rather than wrong, so the repair layer had nothing to
repair. Their causes are the generation budget and the completion path, not the
constraint algorithm.

\paragraph{Compute budget.}
$L$ (span cap) and $k$ (candidates per state) are compute budgets, not tuned
constants; we did not sweep them. With automaton candidates $k$ no longer
controls whether an answer is found, only how many legal alternatives are
weighed, so the search is not silently truncated by the choice.

\section{Conclusion}

Grammar-constrained decoding for diffusion models is a repair problem, and repair
is a search over the parser's own transition structure. Drawing candidates from
that structure rather than from the model's vocabulary ranking removes repair
failure as a category; solving the violated span jointly rather than one token at
a time removes the degeneracy that single-token repair induces; and deriving
termination from the grammar rather than from a step count makes the same rule
work on JSON and SMILES. The resulting system reaches $96.7\%$ schema validity
with $1.2$ repair events per instance, and---unlike the baselines---does not buy
validity by emitting less.

\bibliographystyle{plainnat}
\bibliography{custom}

\end{document}
